% Volume V: Dynamics, Signals, Graphs, and Control
% Continuity note for Chapter 28.
% Previous dependencies: Chapter 12 introduced Fourier analysis and convolution at an operator level; Chapter 15 introduced expectations; Chapter 27 introduced stochastic processes and autocovariance as temporal dependence.
% New durable concepts: sampled time series; sampling interval and sampling rate; continuous traces versus event trains; windows, trials, channels, stationarity; autocorrelation and cross-correlation; discrete Fourier transform; amplitude, phase, power, and spectral leakage; Nyquist frequency and aliasing; linear time-invariant filters; impulse response and frequency response; causal and noncausal filters; neural signal modalities and preprocessing choices.
% Forward hooks: Chapter 29 generalizes signal ideas to graph signals; Chapter 46 uses binned and filtered spike trains; Chapter 48 uses low-dimensional neural trajectories; Chapter 51 uses spectral and time-series summaries for neural-data analysis.
% Notation commitments: Discrete time series use x_0,\ldots,x_{T-1}; sampling interval is \Delta t and sampling rate is f_s=1/\Delta t; DFT coefficients are X_k; autocovariance is C_x(\ell); filters use y_t=\sum_k h_k x_{t-k}.
\section{Chapter 28: Signals, Time Series, and Spectral Analysis}
\label{ch:signals-time-series-spectral-analysis}

Many scientific measurements are functions of time. A voltage trace, EEG channel, local field potential, calcium fluorescence trace, spike-count sequence, behavioral trajectory, training loss curve, and transformer token stream all carry order. Treating such data as an unordered table throws away the structure that often matters most.

\paragraph{Chapter problem.}
How do we represent time-dependent data, measure temporal dependence, decompose signals into frequencies, and design filters without confusing mathematical transformations with biological interpretation?

\paragraph{Section map.}
Section 28.1 defines time-series representation, sampling, trials, channels, event trains, and autocorrelation. Section 28.2 develops spectral analysis through the DFT, amplitude, phase, power, sampling, and aliasing. Section 28.3 defines linear filters, convolution, impulse responses, frequency responses, and filtering failure modes. Section 28.4 specializes the machinery to neural signals and common preprocessing choices.

\subsection{28.1 Time-series representation}

\paragraph{Guiding question.}
What mathematical object represents a time-dependent recording, and what information is carried by its indices and units?

\paragraph{Beginner-facing intuition.}
A time series is not just a vector of numbers. It is a vector whose entries are ordered by time and whose spacing has units. If a signal is sampled every \(\Delta t=1\) ms, then entry \(x_{50}\) refers to time \(50\) ms after the chosen origin. If the same vector is sampled every \(10\) ms, the biological and spectral interpretation changes.

The representation also depends on the measurement type. A local field potential is a continuous-valued trace after sampling. A spike train begins as event times and may be converted into binned counts. A behavioral trajectory may be multivariate, such as hand position and velocity. A neural recording may add trial and channel axes.

\paragraph{Formal definitions and notation.}
A single-channel discrete time series of length \(T\) is
\[
x=(x_0,\ldots,x_{T-1})\in\mathbb R^T.
\]
With sampling interval \(\Delta t>0\), the time of sample \(t\) is
\[
\tau_t=t\Delta t,\qquad t=0,\ldots,T-1.
\]
The sampling rate is
\[
f_s=\frac{1}{\Delta t}.
\]
A multichannel recording can be written
\[
X\in\mathbb R^{T\times C},
\]
where \(X_{t,c}\) is the value at time index \(t\) and channel \(c\). Repeated trials can be stored as
\[
R\in\mathbb R^{B\times T\times C},
\]
where \(B\) is the number of trials.

For a weakly stationary scalar process with mean \(\mu\), the autocovariance at lag \(\ell\) is
\[
C_x(\ell)=\mathbb E[(X_t-\mu)(X_{t+\ell}-\mu)].
\]
For a finite observed signal, a common empirical estimate is
\[
\widehat C_x(\ell)=\frac{1}{T-\ell}\sum_{t=0}^{T-\ell-1}(x_t-\bar x)(x_{t+\ell}-\bar x),
\qquad \ell\geq0.
\]

\paragraph{Core mechanism: binning event times into a time series.}
Suppose spike times in seconds are
\[
0.012,\quad 0.027,\quad 0.031,\quad 0.088.
\]
Using bins of width \(\Delta t=0.02\) s over \([0,0.10)\), the bins are
\[
[0,0.02), [0.02,0.04), [0.04,0.06), [0.06,0.08), [0.08,0.10).
\]
The binned count time series is
\[
x=(1,2,0,0,1).
\]
The same event data can also be represented as a point process
\[
s(t)=\sum_k \delta(t-t_k)
\]
or as a smoothed rate after convolution with a kernel. The representation should match the question: precise timing, count variability, or slowly varying rate.

\paragraph{Concrete example.}
Let
\[
x=(2,4,4,6).
\]
The sample mean is \(\bar x=4\). The lag-one empirical autocovariance estimate with denominator \(T-1=3\) is
\[
\widehat C_x(1)=\frac{1}{3}\bigl[(2-4)(4-4)+(4-4)(4-4)+(4-4)(6-4)\bigr]=0.
\]
The zero-lag estimate is
\[
\widehat C_x(0)=\frac14\left[(-2)^2+0^2+0^2+2^2\right]=2.
\]
This small example shows that autocovariance is not just variance; it asks whether shifted values co-vary.

\paragraph{Computational neuroscience example.}
For an experiment with \(B=40\) trials, \(T=500\) time bins, and \(C=96\) neurons, the binned response tensor
\[
R\in\mathbb R^{40\times500\times96}
\]
has a trial axis, a time axis, and a neuron axis. Averaging over trials estimates a peri-stimulus time course. Averaging over neurons estimates population activity. Computing autocorrelation along the time axis asks whether activity at one bin predicts later activity.

\paragraph{AI or machine-learning example.}
A sequence model receives ordered tokens \(x_1,\ldots,x_T\), not an unordered bag. A training-loss trace \(L_0,\ldots,L_T\) is also a time series: smoothing it, estimating autocorrelation, or checking low-frequency drift can reveal optimization instability or learning-rate effects.

\paragraph{Common misconceptions and failure modes.}
The index \(t\) is not automatically physical time; physical time requires the sampling interval. A time series vector and a feature vector can have the same shape but different meaning. Trial averaging can hide single-trial temporal variability. Autocorrelation estimates assume some stability over the analyzed window; slow drift can masquerade as long memory.

\paragraph{Exercises.}
\begin{enumerate}[label=\textbf{28.1.\arabic*.}, leftmargin=*]
\item \textbf{Mechanical.} A signal is sampled at \(500\) Hz for \(2\) s. What are \(\Delta t\), \(T\), and the time of sample \(t=100\)?
\item \textbf{Proof or derivation.} Show that \(C_x(0)=\operatorname{Var}(X_t)\) for a stationary process.
\item \textbf{Numerical/computational.} Bin spike times \(0.01,0.04,0.041,0.09\) into bins of width \(0.05\) s over \([0,0.10)\).
\item \textbf{Interpretation.} Explain the difference between an EEG time series, an LFP time series, and a binned spike-count time series in terms of units and measurement scale.
\item \textbf{Misconception/debugging.} A student shuffles the time indices of a recording and says the mean and variance are unchanged, so the signal is unchanged. What temporal information has been destroyed?
\end{enumerate}

\subsection{28.2 Spectral analysis}

\paragraph{Guiding question.}
How can a time series be described by oscillatory components, and what do amplitude, phase, and power mean?

\paragraph{Beginner-facing intuition.}
Spectral analysis asks how much of a signal varies slowly, how much oscillates rapidly, and at which frequencies rhythmic structure is concentrated. Instead of asking only ``what is the value at time \(t\)?'', we ask ``which sinusoidal patterns help reconstruct the signal?''

A spectrum is not a magic detector of brain rhythms. It is a transformation of a time series. Peaks can reflect oscillations, artifacts, sharp transients, filtering, or finite-window effects. Interpretation requires the measurement context.

\paragraph{Formal definitions and notation.}
For a length-\(T\) discrete signal \(x_0,\ldots,x_{T-1}\), the discrete Fourier transform is
\[
\boxed{
X_k=\sum_{t=0}^{T-1}x_t\exp\!\left(-2\pi \mathrm{i}\frac{kt}{T}\right),
\qquad k=0,\ldots,T-1.
}
\]
The inverse transform is
\[
x_t=\frac1T\sum_{k=0}^{T-1}X_k\exp\!\left(2\pi \mathrm{i}\frac{kt}{T}\right).
\]
The coefficient \(X_k\) is complex. Its magnitude \(|X_k|\) measures amplitude at the corresponding discrete frequency, and its argument \(\arg(X_k)\) measures phase. A simple periodogram power estimate is proportional to
\[
|X_k|^2.
\]
If the sampling rate is \(f_s\), the physical frequency associated with \(k\) is
\[
f_k=\frac{k}{T}f_s
\]
for low positive frequencies. The Nyquist frequency is
\[
f_{\mathrm{Nyq}}=\frac{f_s}{2}.
\]
Frequencies above Nyquist cannot be represented uniquely after sampling and can alias into lower frequencies.

\paragraph{Core mechanism: orthogonal sinusoidal basis.}
The DFT works because the complex exponentials
\[
\phi_k(t)=\exp\!\left(2\pi \mathrm{i}\frac{kt}{T}\right)
\]
are orthogonal over \(t=0,\ldots,T-1\):
\[
\sum_{t=0}^{T-1}\phi_k(t)\overline{\phi_m(t)}
=
\begin{cases}
T, & k=m,\\
0, & k\neq m.
\end{cases}
\]
Thus \(X_k\) is an inner product of the signal with the \(k\)-th oscillatory basis vector. Spectral analysis is projection onto sinusoidal directions, not a separate kind of data.

\paragraph{Concrete example.}
Let
\[
x=(1,0,-1,0).
\]
For \(T=4\), the DFT coefficient at \(k=1\) is
\[
X_1=1+0\cdot e^{-\pi\mathrm{i}/2}+(-1)e^{-\pi\mathrm{i}}+0\cdot e^{-3\pi\mathrm{i}/2}=1+1=2.
\]
The coefficient at \(k=2\) is
\[
X_2=1+0+(-1)e^{-2\pi\mathrm{i}}+0=0.
\]
The signal is aligned with one cycle over four samples, not with the alternating high-low pattern represented by \(k=2\).

\paragraph{Computational neuroscience example.}
In EEG or LFP analysis, a spectral peak near \(10\) Hz may be called alpha-band power if it appears in the relevant recording context. The spectrum alone does not prove a mechanism. A transient artifact, a sharp evoked response, or volume conduction can affect power. Good analysis reports sampling rate, preprocessing, window length, reference choice, and whether the spectrum is trial-averaged or estimated from continuous data.

\paragraph{AI or machine-learning example.}
Sequence models and convolutional networks can be probed spectrally. A learned temporal filter may suppress high-frequency noise while retaining slower trends. In audio and speech models, spectrograms represent local frequency content over sliding windows. In optimization, oscillatory loss components can reveal unstable learning rates or momentum-induced ringing.

\paragraph{Common misconceptions and failure modes.}
A power spectrum loses phase information unless phase is stored separately. A spectral peak is not automatically a sustained oscillation; finite windows and sharp transients spread power across frequencies. Windowing reduces edge artifacts but changes the estimator. Sampling too slowly causes aliasing, where a high-frequency component appears as a lower-frequency one.

\paragraph{Exercises.}
\begin{enumerate}[label=\textbf{28.2.\arabic*.}, leftmargin=*]
\item \textbf{Mechanical.} A signal sampled at \(1000\) Hz has Nyquist frequency \(500\) Hz. What is the Nyquist frequency if the sampling rate is \(200\) Hz?
\item \textbf{Proof or derivation.} Prove the DFT orthogonality identity for \(k\neq m\) by summing a finite geometric series.
\item \textbf{Numerical/computational.} Compute \(X_0\) and \(X_2\) for \(x=(1,-1,1,-1)\).
\item \textbf{Interpretation.} What is lost if you keep only \(|X_k|^2\) and discard the phases?
\item \textbf{Misconception/debugging.} A student sees a \(60\) Hz peak in an LFP spectrum and calls it a neural oscillation. What artifact or measurement issue should be checked first?
\end{enumerate}

\subsection{28.3 Linear filters}

\paragraph{Guiding question.}
How does a filter transform a time series, and why does convolution become multiplication in the frequency domain?

\paragraph{Beginner-facing intuition.}
A filter is a rule that takes an input signal and produces an output signal. A moving average smooths rapid fluctuations. A high-pass filter removes slow drift. A band-pass filter keeps a chosen frequency band. The simplest filters are linear and time-invariant: doubling the input doubles the output, adding inputs adds outputs, and shifting the input shifts the output.

Filtering is useful, but it is also dangerous. A filter can delay a signal, smear events, remove meaningful activity, or create ringing. The mathematical design should be tied to a measurement question.

\paragraph{Formal definitions and notation.}
A discrete-time filter is \emph{linear} if
\[
F(ax+bz)=aF(x)+bF(z)
\]
for signals \(x,z\) and scalars \(a,b\). It is \emph{time-invariant} if shifting the input shifts the output by the same amount. A linear time-invariant filter is determined by its impulse response \(h_k\), and its output is convolution:
\[
\boxed{y_t=(h*x)_t=\sum_k h_k x_{t-k}.}
\]
A finite impulse response filter has only finitely many nonzero \(h_k\). A causal filter has \(h_k=0\) for \(k<0\), so it uses only present and past inputs.

The frequency response is the Fourier transform of the impulse response:
\[
H(\omega)=\sum_k h_k e^{-\mathrm{i}\omega k}.
\]
For an input sinusoid \(x_t=e^{\mathrm{i}\omega t}\), the output is
\[
y_t=H(\omega)e^{\mathrm{i}\omega t}.
\]
Thus \(H(\omega)\) tells us the gain and phase shift applied to frequency \(\omega\).

\paragraph{Core mechanism: convolution theorem.}
For finite circular convolution, if
\[
y_t=\sum_{k=0}^{T-1}h_k x_{t-k \ \mathrm{mod}\ T},
\]
then the DFT satisfies
\[
\boxed{Y_m=H_mX_m.}
\]
\emph{Derivation.} Substitute the convolution into the DFT:
\[
Y_m=\sum_t\sum_k h_k x_{t-k}e^{-2\pi\mathrm{i}mt/T}.
\]
Let \(u=t-k\). Then
\[
Y_m=\sum_k h_ke^{-2\pi\mathrm{i}mk/T}\sum_u x_u e^{-2\pi\mathrm{i}mu/T}
=H_mX_m.
\]
Filtering in time by convolution becomes pointwise multiplication in frequency.

\paragraph{Concrete example.}
The three-point causal moving average has
\[
h_0=h_1=h_2=\frac13.
\]
For \(x=(3,6,0,0)\), using zeros before the signal begins,
\[
y_0=1,\qquad y_1=3,\qquad y_2=3,\qquad y_3=2.
\]
The output is smoother but delayed: the large value at \(x_1=6\) affects \(y_1,y_2,y_3\).

\paragraph{Computational neuroscience example.}
Smoothing binned spikes with a Gaussian or exponential kernel estimates an instantaneous firing-rate trace. A causal exponential filter respects time direction and resembles synaptic filtering. A noncausal centered Gaussian is useful for offline visualization but uses future spikes to estimate the current rate, so it should not be used for real-time decoding without acknowledging the look-ahead.

\paragraph{AI or machine-learning example.}
One-dimensional convolutional layers in sequence models are learned filters. Their impulse responses are learned parameters, and their frequency responses reveal which temporal scales they emphasize. In signal preprocessing for speech or neural decoding, fixed filters may remove low-frequency drift or high-frequency noise before a learned model sees the data.

\paragraph{Common misconceptions and failure modes.}
Smoothing is not the same as denoising; it can remove true fast signals. Zero-phase offline filtering is noncausal and can leak future information into past estimates. Sharp frequency cutoffs can cause ringing in time. Filtered data should not be interpreted as raw data; the filter has changed amplitudes, phases, and temporal resolution.

\paragraph{Exercises.}
\begin{enumerate}[label=\textbf{28.3.\arabic*.}, leftmargin=*]
\item \textbf{Mechanical.} For \(h=(1/2,1/2)\) and \(x=(2,4,8)\), compute the causal filtered output using zeros before \(t=0\).
\item \textbf{Proof or derivation.} Show that if \(x_t=e^{\mathrm{i}\omega t}\), then convolution by \(h\) produces \(y_t=H(\omega)e^{\mathrm{i}\omega t}\).
\item \textbf{Numerical/computational.} Compare a causal moving average \(h_0=h_1=h_2=1/3\) with a centered noncausal average on the impulse signal \(x=(0,0,1,0,0)\).
\item \textbf{Interpretation.} Why might a filter that removes low frequencies also remove meaningful slow neural state changes?
\item \textbf{Misconception/debugging.} A student filters test data using a centered window and then evaluates real-time decoding performance. What information leak may have occurred?
\end{enumerate}

\subsection{28.4 Neural signals}

\paragraph{Guiding question.}
How do signal-processing choices change what can be inferred from neural recordings?

\paragraph{Beginner-facing intuition.}
Neural data are not one kind of signal. A spike train records discrete events. Membrane voltage records a cellular state. LFP and EEG summarize extracellular fields over space and time. Calcium imaging is slow and indirect relative to spikes. Functional MRI is much slower still and reflects hemodynamic changes. Each modality has its own sampling scale, units, noise sources, and preprocessing hazards.

Mathematics helps by forcing explicit choices: what is the signal, what are the axes, what is the sampling rate, what transformations were applied, and what biological claim is supported by those transformed data?

\paragraph{Formal definitions and notation.}
A spike train with spike times \(\{t_k\}\) can be represented as
\[
s(t)=\sum_k \delta(t-t_k).
\]
Binning with width \(\Delta\) gives counts
\[
n_j=\int_{j\Delta}^{(j+1)\Delta}s(t)\,dt.
\]
Smoothing with kernel \(g\) gives a rate estimate
\[
\hat r(t)=(g*s)(t)=\sum_k g(t-t_k).
\]
For continuous neural traces \(v_t\), preprocessing may include detrending, line-noise removal, band-pass filtering, trial alignment, artifact rejection, and normalization. Each step is a mathematical transformation, not merely a cosmetic operation.

\begin{center}
\small
\renewcommand{\arraystretch}{1.18}
\begin{tabular}{p{0.18\linewidth}p{0.18\linewidth}p{0.24\linewidth}p{0.30\linewidth}}
\toprule
Signal type & Typical time scale & Mathematical representation & Common caution \\
\midrule
Spikes & milliseconds & event times, point process, binned counts & bin width and smoothing change timing information \\
Membrane voltage & sub-millisecond to milliseconds & sampled continuous trace & electrode noise and filtering affect thresholds \\
LFP/ECoG/EEG & milliseconds to seconds & multichannel continuous time series & reference, volume conduction, and artifacts matter \\
Calcium imaging & tens to hundreds of milliseconds & fluorescence time series & slow indicator dynamics blur spikes \\
fMRI & seconds & voxel or region time series & hemodynamic response is indirect and slow \\
\bottomrule
\end{tabular}
\end{center}

\paragraph{Core mechanism: from spikes to smoothed rate.}
Let spike times be \(0.01,0.04,0.09\) s and let bins have width \(0.05\) s on \([0,0.10)\). The binned counts are
\[
n=(2,1).
\]
The binned rate estimate is
\[
\hat r=(2/0.05,\ 1/0.05)=(40,\ 20)\ \text{spikes/s}.
\]
If instead we smooth with an exponential kernel
\[
g(t)=\frac{1}{\tau}e^{-t/\tau}\mathbf 1_{t\geq0},
\]
then each spike contributes a causal pulse that decays with time constant \(\tau\). The binned estimate emphasizes window counts; the smoothed estimate emphasizes a continuous causal trace.

\paragraph{Concrete example.}
Suppose an LFP trace is sampled at \(1000\) Hz and analyzed in \(1\) s windows. The frequency resolution of a basic DFT periodogram is \(1\) Hz. If the window is shortened to \(0.25\) s, resolution becomes \(4\) Hz. Short windows localize time better but blur frequency more. This time-frequency tradeoff is a mathematical fact, not a neuroscience-specific inconvenience.

\paragraph{Computational neuroscience example.}
In a motor-cortex task, spikes may be binned at \(20\) ms for decoding hand velocity, while LFP beta-band power may be estimated over longer windows. The spike bins preserve rapid population changes for decoding. The spectral power estimate summarizes rhythmic activity over a window. Combining them requires acknowledging that the features have different temporal resolution and different biological sources.

\paragraph{AI or machine-learning example.}
Neural time-series pipelines often feed filtered, binned, or spectral features into classifiers, recurrent networks, or state-space models. A model trained on preprocessed signals learns from the transformed representation. If preprocessing leaks future information, changes label timing, or removes task-relevant bands, the downstream ML result can be misleading even if the model architecture is strong.

\paragraph{Common misconceptions and failure modes.}
No preprocessing step is neutral. A spike rate is an estimate, not the raw spike train. Band power is not a direct measurement of a single oscillatory generator. Trial alignment can create sharp averages even when single trials vary in latency. High decoding accuracy from filtered data does not by itself identify the neural mechanism used by the brain.

\paragraph{Exercises.}
\begin{enumerate}[label=\textbf{28.4.\arabic*.}, leftmargin=*]
\item \textbf{Mechanical.} Convert binned spike counts \((0,3,1)\) in \(50\) ms bins into firing rates in spikes/s.
\item \textbf{Proof or derivation.} Show that binning a point process by integration gives integer counts when no spike lies exactly on a bin boundary.
\item \textbf{Numerical/computational.} An LFP signal sampled at \(500\) Hz is analyzed in \(2\) s windows. What is the number of samples per window and the basic frequency spacing?
\item \textbf{Interpretation.} Why should calcium imaging traces not be interpreted as instantaneous spike trains?
\item \textbf{Misconception/debugging.} A student reports ``theta power increased,'' but does not specify sampling rate, filter band, window length, or artifact rejection. What information is missing for interpreting the claim?
\end{enumerate}

\paragraph{Closing bridge.}
Signals carry structure along a time axis. We represented samples, events, spectra, filters, and neural measurement pipelines with explicit units and transformations. The next chapter keeps the same discipline but replaces the time axis with a network: signals can live on graph nodes, and graph Laplacian eigenvectors play a role analogous to Fourier modes.

\clearpage
